\chapter{$Q$, the Running Average (and What $x_n$ Physically Is)}
\label{ch:q_average}

\begin{studentq}
``So the $x_n$ are the values from the value heads, that is $V$'s of each position arising when a new candidate is visited. Correct?'' --- Yes, exactly.
\end{studentq}

Each time the search goes down arrow $a$ and eventually evaluates some position at the bottom, one number comes back. Call the $n$-th such number $x_n$.

\begin{align*}
  W(a) &= x_1 + x_2 + \dots + x_n \quad \text{(running total)} \\
  n(a) &= \text{how many times we went down this arrow} \\
  Q(a) &= \frac{W(a)}{n(a)} \quad \text{(the average)}
\end{align*}

$Q$ is a \textbf{sample mean}. Nothing more exotic than that.

The incremental form (which is what the code actually does, so it never has to store the list) is:
\begin{equation}
  Q_{\text{new}} = Q_{\text{old}} + \frac{x_n - Q_{\text{old}}}{n}
\end{equation}

Check it once by hand, with $x_1 = 0.968$ then $x_2 = 0.900$:
\begin{align*}
  n=1: \quad & Q = 0.968 \\
  n=2: \quad & Q = 0.968 + \frac{0.900 - 0.968}{2} = 0.968 - 0.034 = \mathbf{0.934} \\
  \text{check:} \quad & \frac{0.968 + 0.900}{2} = \mathbf{0.934} \quad \text{(agrees)}
\end{align*}

\begin{center}
\begin{minipage}{\linewidth}
\centering
% fig_c_qmean.tex -- Sample mean convergence and step damping
\begin{tikzpicture}[scale=0.88, every node/.style={font=\small}]
  % Axes
  \draw[WorkGrey, line width=0.7pt, ->] (0,0) -- (7.5, 0) node[right, font=\footnotesize] {Visit count $n$};
  \draw[WorkGrey, line width=0.7pt, ->] (0,0) -- (0, 4.2) node[above, font=\footnotesize] {Evaluation ($x_n, Q_n$)};

  % Y-ticks
  \foreach \y/\val in {0.8/0.80, 2.0/0.90, 3.2/0.95, 4.0/1.00} {
    \draw[WorkGrey!30, line width=0.3pt] (0, \y) -- (7.0, \y);
    \node[left, font=\tiny, text=WorkGrey] at (0, \y) {\val};
  }

  % X-ticks
  \foreach \x in {1, 2, 3, 4, 5, 6} {
    \draw[WorkGrey, line width=0.5pt] (\x, 0.05) -- (\x, -0.05);
    \node[below, font=\footnotesize] at (\x, -0.05) {$n{=}\x$};
  }

  % Raw samples x_n
  \node[draw=PitfallRed, fill=PitfallRed!20, star, star points=5, star point ratio=2, inner sep=1.5pt] (x1) at (1, 3.56) {};
  \node[draw=PitfallRed, fill=PitfallRed!20, star, star points=5, star point ratio=2, inner sep=1.5pt] (x2) at (2, 2.00) {};
  \node[draw=PitfallRed, fill=PitfallRed!20, star, star points=5, star point ratio=2, inner sep=1.5pt] (x3) at (3, 3.40) {};
  \node[draw=PitfallRed, fill=PitfallRed!20, star, star points=5, star point ratio=2, inner sep=1.5pt] (x4) at (4, 2.24) {};
  \node[draw=PitfallRed, fill=PitfallRed!20, star, star points=5, star point ratio=2, inner sep=1.5pt] (x5) at (5, 3.32) {};
  \node[draw=PitfallRed, fill=PitfallRed!20, star, star points=5, star point ratio=2, inner sep=1.5pt] (x6) at (6, 2.84) {};

  % Running mean line Q_n
  \draw[BuildBlue, line width=1.5pt] 
    (1, 3.56) -- (2, 2.82) -- (3, 3.03) -- (4, 2.82) -- (5, 2.94) -- (6, 2.91);

  \foreach \x/\y/\lbl in {1/3.56/0.968, 2/2.82/0.934, 3/3.03/0.943, 4/2.82/0.934, 5/2.94/0.939, 6/2.91/0.938} {
    \fill[BuildBlue] (\x, \y) circle (2.5pt);
    \node[above=2pt, font=\tiny\bfseries, text=NavyBlue] at (\x, \y) {\lbl};
  }

  % Annotations
  \draw[GoldPath, line width=1.0pt, {Stealth[length=2mm]}-] (2, 2.7) -- ++(0.6, -0.6)
    node[right, font=\tiny, align=left, fill=white, inner sep=1pt] {Step size $\frac{1}{2}$:\\$Q$ shifts by $-0.034$};
    
  \draw[GoldPath, line width=1.0pt, {Stealth[length=2mm]}-] (6, 2.95) -- ++(0.6, 0.6)
    node[right, font=\tiny, align=left, fill=white, inner sep=1pt] {Step size $\frac{1}{6}$:\\$Q$ barely moves};

  % Legend
  \node[anchor=north west, draw=WorkGrey!40, fill=white, rounded corners=2pt, font=\scriptsize] at (0.3, 1.4) {%
    \begin{tabular}{cl}
      \textcolor{PitfallRed}{$\star$} & Raw leaf sample $x_n$ (volatile)\\[0.3ex]
      \textcolor{BuildBlue}{$\bullet$---$\bullet$} & Running mean $Q_n = Q_{n-1} + \frac{1}{n}(x_n - Q_{n-1})$ (convergent)
    \end{tabular}%
  };
\end{tikzpicture}

\captionof{figure}{\textbf{Running-mean convergence and step damping.} Volatile leaf evaluations $x_n$ are smoothed into a stable sample mean $Q_n$ via step size $1/n$.}
\label{fig:qmean}
\end{minipage}
\end{center}

\section{The Sign Flip, Carefully}

$x_n$ arrives from a position where it is the \textit{opponent's} turn. The network reports $V$ from the point of view of whoever is to move there. So before adding it into our total we negate it:
\begin{equation}
  x_n = -V(\text{child position})
\end{equation}

In our endgame, after \texttt{Kd6} it is Black to move, and the network says Black stands at about \textbf{$-0.968$} (Black is losing). White's edge records:
\begin{equation}
  x_1 = -(-0.968) = \mathbf{+0.968}
\end{equation}

\begin{center}
\begin{minipage}{\linewidth}
\centering
% fig_c_signflip.tex -- Negamax sign flip across alternating plies
\begin{tikzpicture}[
    levelnode/.style={rounded corners=3pt, draw=NavyBlue, fill=PaleGrey,
                      minimum width=46mm, minimum height=11mm, align=center, font=\scriptsize},
    downedge/.style={-{Stealth[length=2.5mm]}, draw=NavyBlue, line width=1.0pt},
    upedge/.style={-{Stealth[length=2.5mm]}, draw=IdeaGreen, line width=1.2pt, dashed}
  ]

  % Level 0: Root (White to move)
  \node[levelnode, fill=BuildBlue!12] (root) at (0, 3.2) {%
    \textbf{Root: White to Move}\\[0.3ex]
    White's evaluation: \textbf{$V = +0.97602$ (Winning!)}\\
    \color{BuildBlue}\textbf{Recorded on Kd6 edge: $Q = \mathbf{+0.96766}$}%
  };

  % Level 1: Child after Kd6 (Black to move)
  \node[levelnode, fill=PitfallRed!10] (child) at (0, 0) {%
    \textbf{Child: Black to Move (after 1. Kd6)}\\[0.3ex]
    Network evaluates from \textbf{Black's POV}:\\
    \color{PitfallRed}\textbf{$V(\text{child}) = \mathbf{-0.96766}$ (Black is losing!)}%
  };

  % Level 2: Grandchild after Kd8 (White to move)
  \node[levelnode, fill=BuildBlue!12] (grandchild) at (0, -3.2) {%
    \textbf{Grandchild: White to Move (after 1... Kd8)}\\[0.3ex]
    Network evaluates from \textbf{White's POV}:\\
    \color{IdeaGreen}\textbf{$V(\text{grandchild}) = \mathbf{+0.95128}$ (White is winning!)}%
  };

  % Forward move edges
  \draw[downedge] (root.south west) ++(1.0, 0) -- 
    node[left, font=\tiny\bfseries, text=NavyBlue] {1. Kd6} 
    ($(child.north west)+(1.0, 0)$);

  \draw[downedge] (child.south west) ++(1.0, 0) -- 
    node[left, font=\tiny\bfseries, text=NavyBlue] {1... Kd8} 
    ($(grandchild.north west)+(1.0, 0)$);

  % Backpropagation Negamax sign-flip arrows
  \draw[upedge] ($(grandchild.north east)+(-1.0, 0)$) -- 
    node[right, font=\tiny\bfseries, text=IdeaGreen, align=left] {Sign flip:\\[0.2ex] $x = -V = -0.95128$} 
    ($(child.south east)+(-1.0, 0)$);

  \draw[upedge] ($(child.north east)+(-1.0, 0)$) -- 
    node[right, font=\tiny\bfseries, text=IdeaGreen, align=left] {Sign flip:\\[0.2ex] $x_1 = -(-0.96766) = \mathbf{+0.96766}$} 
    ($(root.south east)+(-1.0, 0)$);

  % Side explanation box
  \node[draw=IdeaGreen, fill=IdeaGreen!8, rounded corners=3pt, font=\scriptsize, align=left, anchor=west] at (3.2, 0) {%
    \textbf{The Negamax Invariant:}\\[0.5ex]
    Because chess is zero-sum,\\
    whatever is good for Black is bad for White.\\[0.8ex]
    $\mathbf{x_n = -V(\text{child position})}$\\[0.8ex]
    A single sign flip handles the alternating\\
    points of view at every tree depth.%
  };
\end{tikzpicture}

\captionof{figure}{\textbf{The Negamax sign-flip across alternating plies.} $x_n = -V(\text{child})$ guarantees that every node accumulates value from its own perspective.}
\label{fig:signflip}
\end{minipage}
\end{center}

This is why the guide's figure shows the value ``flipping sign at every level.'' It is the same fact as $V_{\text{parent}} = -V_{\text{child}}$ from Chapter~\ref{ch:network_outputs}.
